\chapter{Probes: reading a mind, and the traps in claiming you did}
\label{ch:probes}

\begin{objectives}
\begin{itemize}
\item Define an activation, and say exactly what object a probe is trained on.
\item Build a linear probe --- and recognise it as the model of Chapter~\ref{ch:learned}, with
  nothing added.
\item Name the four ways a probe result can be worthless, and the control that defuses each.
\item Distinguish \emph{information is present} from \emph{information is used}, and describe
  the intervention that separates them.
\item Write down what a probe with $92\%$ accuracy does and does not entitle you to say.
\end{itemize}
\end{objectives}

\section{What there is to read}

A forward pass through BT3 produces, at every layer, a $64\times768$ block of numbers: for
each square, 768 quantities the network computed on its way to an answer. These are
\textbf{activations}. They are not designed to mean anything to us --- nobody labelled
dimension 412 --- but they are the only substrate the network has, so if the network
"understands" anything, that understanding is a pattern in these numbers.

\begin{keyidea}
Interpretability is not mind-reading. It is a much humbler claim: \textbf{if a property of the
position can be recovered from the activations by a simple rule, then the activations encode
that property.} Everything hard about the subject is in the words "simple" and "encode".
\end{keyidea}

\section{A probe is Chapter 10, applied to the middle of a network}

Fix a layer $l$ and a property you care about --- say, \emph{is square e5 attacked by White?}
Collect a large set of positions, run each through the network, and record two things: the
activation vector $\vv{x}^{(l)}_i \in \R^{768}$ at the square of interest, and the true
answer $y \in \{0,1\}$, which you compute with an ordinary chess library. Now fit

\begin{equation}
  \label{eq:probe}
  \hat{y} \;=\; \sigma\bigl(\vv{w}\cdot\vv{x}^{(l)}_i + b\bigr)
\end{equation}

by the gradient descent of Chapter~\ref{ch:learned}, with cross-entropy loss. That is the
entire method. A \textbf{linear probe} is logistic regression on frozen activations.

\begin{notationbox}
\emph{Frozen} means the network's weights are not updated --- only $\vv{w}$ and $b$ are
learned. The network is a fixed function; the probe is a ruler held against its insides.
\end{notationbox}

Two design choices carry all the meaning:

\begin{itemize}
\item \textbf{Which layer.} Probing layer 1 tells you about the input encoding; probing layer
  14 tells you about what is about to become the output. The interesting claims are always
  about the middle.
\item \textbf{How simple the probe is.} This is the crux, and §13.3 is about why.
\end{itemize}

\begin{worked}{what a probe experiment looks like end to end}
\emph{A concrete protocol, written out so it can be executed. Nothing here has been run in
this project --- see the confidence note below.}
\begin{enumerate}
\item Sample 50{,}000 positions from the corpus in \texttt{games\_of\_derdiedasdie/}.
\item For each, compute the label with \texttt{python-chess}: \emph{is the black king in
  check within 2 plies against best play?} (or whatever property is under test).
\item Run each position through the net with hooks on layer $l$ (\texttt{lczerolens}),
  recording the $64\times768$ activation block.
\item Split 80/20 into train and test, by \emph{game}, not by position --- positions from the
  same game are near-duplicates and will leak.
\item Fit Equation~\eqref{eq:probe}; report accuracy on the held-out games.
\item Report the three baselines of §13.3 alongside it. Without them the number means nothing.
\end{enumerate}
\end{worked}

\section{The four ways to fool yourself}

A probe reaching $92\%$ accuracy feels like a discovery. It is not, on its own. Here are the
four failures, each with the control that defuses it.

\subsection{Trap 1: the probe learned the task itself}

If a property is easy to compute from the board, a probe may be computing it from the
activations' \emph{input} content rather than reading a representation the network built.
A probe on layer 1 --- which is nearly the raw board --- can detect "there is a rook on this
square" perfectly, and that tells you nothing about the network's understanding.

\begin{keyidea}
\textbf{Control 1: the input baseline.} Train the identical probe on the \emph{input}
encoding $X^{(0)}$. The claim "layer 8 encodes property $p$" is only interesting to the extent
that layer 8 beats layer 0. Report both numbers, always. A layer-8 accuracy of $92\%$ against
a layer-0 accuracy of $90\%$ is a non-result dressed as a result.
\end{keyidea}

\subsection{Trap 2: a strong probe can extract anything}

Give the probe a hidden layer, or two, and it becomes a neural network in its own right ---
capable of \emph{computing} the property from raw material rather than reading it off. Taken
far enough, "the activations encode $p$" degenerates into "the activations contain the
position, and $p$ is a function of the position", which is true and empty.

\begin{keyidea}
\textbf{Control 2: keep the probe linear, and report its capacity.} The standard convention
is a single linear layer for exactly this reason: if a \emph{linear} rule recovers the
property, the network has done the hard part and laid the answer out in a readable way. If
you need a two-layer probe, say so, and weaken the claim accordingly.
\end{keyidea}

\subsection{Trap 3: present but unused}

This is the deepest one. Suppose layer 8 does encode "there is a knight fork available". It
does not follow that the network \emph{uses} that information when choosing its move. Networks
compute a great deal that ends up discarded --- residual streams carry information forward
that later layers ignore.

The distinction is not academic for your project: a suppressed-win detector is precisely a
claim about information that is present and \emph{not} used, so the two must be measurable
separately or the claim is unfalsifiable.

\begin{keyidea}
\textbf{Control 3: intervene.} Correlation is established by probing; causation requires
changing something and seeing the output move. Two standard interventions:
\begin{itemize}
\item \textbf{Ablation:} zero out (or randomise) the activations at layer $l$ on the squares
  in question, rerun the remaining layers, and see whether the policy changes. If it does not,
  the information was not being used.
\item \textbf{Activation patching:} take the activations from a \emph{different} position ---
  identical except that the property is false --- and splice them in. If the output flips to
  what the other position's output would be, you have shown that those activations carry the
  effect.
\end{itemize}
\end{keyidea}

\subsection{Trap 4: you searched a large space and reported the maximum}

BT3 has 15 layers and 24 heads. Trying every combination and reporting the best is a
multiple-comparisons problem: with 360 candidates, something will look striking by chance.
Anyone who has run a backtest knows this failure by another name.

\begin{keyidea}
\textbf{Control 4: pre-register and re-test.} Decide the layer and the property \emph{before}
looking, or --- if you must search --- hold out a completely separate set of positions and
re-run only the winner on it. Report how many candidates you tried. A layer chosen from 360 by
its performance needs its performance re-measured somewhere clean.
\end{keyidea}

\section{A worked interpretation}

Suppose all four controls are run and the result is:

\begin{center}
\begin{tabular}{@{}lr@{}}
\toprule
Linear probe on layer 8 & $92\%$ \\
Linear probe on layer 0 (input baseline) & $61\%$ \\
Linear probe on a randomly-initialised network, layer 8 & $58\%$ \\
Majority-class baseline & $55\%$ \\
Ablating the relevant activations changes the top policy move & in $40\%$ of cases \\
\bottomrule
\end{tabular}
\end{center}
\emph{(Illustrative numbers, for the purpose of showing what a full report looks like.)}

\begin{center}
\begin{tabular}{@{}p{6.6cm}p{6.6cm}@{}}
\toprule
\textbf{You may say} & \textbf{You may not say} \\
\midrule
"A linear function of layer-8 activations predicts $p$ at $92\%$, against a $61\%$ input
baseline." & "The network knows $p$." \\
\addlinespace
"The information is linearly available at layer 8 and largely absent at the input." &
"Layer 8 computes $p$." (You have shown availability, not the mechanism.) \\
\addlinespace
"Ablating it changes the chosen move in $40\%$ of positions, so it is used at least
sometimes." & "The network's move is caused by $p$." \\
\addlinespace
"On this network, this property, this position distribution." & "Chess networks represent
$p$." \\
\bottomrule
\end{tabular}
\end{center}

\begin{pitfall}
The right-hand column is not pedantry --- it is the difference between a claim your tool can
show a user and a claim that will eventually be embarrassing. This project's own history
contains a metric that called quiet moves "sacrifices" because nothing checked material, and a
batch of fabricated master annotations. Both were plausible, fluent and wrong. An
interpretability result is far easier to produce in that style than a chess fact is, because
nobody can check it by looking at the board.
\end{pitfall}

\section{What probes are actually good for here}

Given all those caveats, why bother? Because the alternative --- reading only the final
outputs --- discards everything the network computed on the way, and the project's north star
is precisely to recover that.

Three uses, in increasing order of ambition and decreasing order of safety:

\begin{enumerate}
\item \textbf{Verification.} Probe for something you can compute anyway (is this piece
  pinned?) purely to establish that the layer is legible at all. Cheap, low-risk, and it
  calibrates your baselines before you go looking for anything exciting.
\item \textbf{Feature discovery.} Probe for properties your fact extractor already emits, to
  find which of them the network actually represents. A fact the network does not encode is a
  fact it cannot be "thinking about", which is directly relevant to the salience problem of
  Chapter~\ref{ch:speak}.
\item \textbf{Suppressed-win detection.} Probe for "a winning tactic exists here" and compare
  with what the policy head does. This is Chapter~\ref{ch:lookahead}, and it is the one that
  needs every control in this chapter.
\end{enumerate}

\begin{hypothesis}
Nothing in this chapter has been run on \texttt{BT3-768x15x24h} in this project. The hooks
exist (\texttt{backend/neural\_vision.py} already extracts intermediate activations for the
attention display), so the machinery is close --- but no probe has been trained, no baseline
measured, and no ablation performed. Until then, every claim of the form "the network
internally sees $X$" is a plan, not a finding.
\end{hypothesis}

\begin{repobox}[title={in this repository}]
The realistic first experiment is use 1 above: probe for a property your relational fact
extractor computes deterministically (an absolute pin, say), on positions from your own
corpus, with the input baseline. It is a day of work, it is falsifiable, and it establishes
whether the interesting experiments are worth attempting on this net.
\end{repobox}

\begin{exercises}

\exhead[the object]{ch13:object} What exactly does a probe take as input --- one number, 768
numbers, or $64\times768$? Describe two different probe designs (per-square and whole-board)
and a property suited to each.

\exhead[probe = chapter 10]{ch13:same} Write down Equation~\eqref{eq:probe} and the model from
§10.2.3 side by side. What differs, apart from where the input comes from? What is frozen in
each case?

\exhead[baselines]{ch13:baselines} A colleague reports "our layer-6 probe detects backward
pawns at $88\%$". List the three numbers you must ask for before believing anything, and say
what each would rule out.

\exhead[capacity]{ch13:capacity} Why is a two-layer probe reaching $95\%$ weaker evidence than
a linear probe reaching $85\%$? Answer in terms of what the probe might be doing rather than
reading.

\exhead[present vs used]{ch13:presentused} Construct a hypothetical case where a property is
perfectly linearly decodable at layer 10 and yet has zero influence on the network's output.
What experiment distinguishes your case from a genuine one?

\exhead[designing an ablation]{ch13:ablation} You want to test whether the network's
representation of a knight fork on e6 is used. Specify: what you zero out, what you rerun,
what you measure, and what result would falsify the claim.

\exhead[patching]{ch13:patching} Activation patching needs a pair of positions differing in
one property. Explain why constructing such pairs is harder in chess than in language, and
propose a way to build 100 such pairs from your corpus.

\exhead[multiple comparisons]{ch13:multcomp} You try 15 layers and report the best. If each
layer's accuracy were pure noise around $60\%$ with a standard deviation of $2$ points, what
would the maximum of 15 draws look like? What does this imply about the headline you would
have written?

\exhead[writing the claim]{ch13:claim} Rewrite this sentence so it would survive review:
"Layer 8 sees the winning sacrifice." Give the version you would put in a paper and the
version you would show a user in the tool.

\exhead[the cheap first experiment]{ch13:first} Design the smallest probe experiment that
would tell you whether probing this net is worth pursuing at all. Specify the property, the
positions, the baseline, the number you would need to see, and the number at which you would
abandon the line of work.

\end{exercises}
